\chapter{Grundlagen}
\label{ch:grundlagen}

\section{Verweise mit cleveref}
\label{sec:verweise}

\Cref{tab:befehle} fasst die wichtigsten Befehle zusammen.
Am Satzanfang steht \verb|\Cref|, im Satz \verb|\cref|.
Mehrere Ziele auf einmal: \cref{ch:einleitung,ch:fazit}.

\begin{table}[ht]
  \centering
  \begin{tabular}{lp{7.5cm}}
    \toprule
    Befehl & Wirkung \\
    \midrule
    \verb|\input{datei}|   & fügt die Datei an dieser Stelle ein \\
    \verb|\include{datei}| & wie \verb|\input|, zusätzlich neue Seite und eigene \texttt{.aux}-Datei \\
    \verb|\includeonly{...}| & übersetzt nur die genannten Dateien \\
    \bottomrule
  \end{tabular}
  \caption{Befehle zum Aufteilen eines Dokuments}
  \label{tab:befehle}
\end{table}

\section{Abkürzungen}

\gls{ci} kam schon in der Einleitung vor und erscheint hier nur noch als Kurzform.
\gls{cd} wird zum ersten Mal verwendet und daher ausgeschrieben.
\todo[inline]{Abschnitt zu Continuous Delivery ausbauen}
